% Generated draft: requires mathematical and source review.
\documentclass[11pt]{article}
\usepackage[margin=25mm]{geometry}
\usepackage{fontspec}
\setmainfont{Latin Modern Roman}
\usepackage{amsmath,amssymb,mathtools,unicode-math}
\setmathfont{Latin Modern Math}
\usepackage{enumitem,needspace,xurl,hyperref}
\setlength{\parindent}{0pt}
\setlength{\parskip}{5pt}
\setlength{\emergencystretch}{3em}
\newcommand{\N}{\mathbb{N}}
\newcommand{\Z}{\mathbb{Z}}
\newcommand{\Q}{\mathbb{Q}}
\newcommand{\R}{\mathbb{R}}
\newcommand{\C}{\mathbb{C}}
\newcommand{\F}{\mathbb{F}}
\newcommand{\A}{\mathbb{A}}
\newcommand{\PP}{\mathbb{P}}
\newcommand{\E}{\mathbb{E}}
\newcommand{\eps}{\varepsilon}
\newcommand{\dd}{\,\mathrm{d}}
\begin{document}
\section*{UnsolvedMath ID: 30002220}
\textbf{Source problem code:} catalogue entry
\subsection{Definitions and mathematical statement}
Let \(A\) be a complex \(C^*\)-algebra, and let \(A_{\mathrm{sa}}=\{a\in A:a=a^*\}\). For an integer \(m\geq 1\), define
\[
\operatorname{Gen}_m(A)_{\mathrm{sa}}=\{(a_1,\ldots,a_m)\in A_{\mathrm{sa}}^m:C^*(a_1,\ldots,a_m)=A\},
\]
where \(C^*(a_1,\ldots,a_m)\) denotes the smallest norm-closed \(*\)-subalgebra containing the listed elements, without adjoining an identity by convention. Give \(A_{\mathrm{sa}}^m\) the maximum norm. For unital \(A\), its generator rank is
\[
\operatorname{gr}(A)=\min\{k\in\{0,1,2,\ldots\}:\operatorname{Gen}_{k+1}(A)_{\mathrm{sa}}\text{ is norm dense in }A_{\mathrm{sa}}^{k+1}\},
\]
with value \(\infty\) if this set is empty.

Separable means having a countable norm-dense subset. Write \(\mathcal Z\) for the Jiang--Su algebra. The algebra \(A\) absorbs \(\mathcal Z\), or is \(\mathcal Z\)-stable, if \(A\cong A\otimes_{\min}\mathcal Z\). Nuclearity means that, for every complex \(C^*\)-algebra \(B\), the minimal and maximal \(C^*\)-tensor norms on \(A\odot B\) coincide. For unital \(A\), real rank zero means that self-adjoint elements with finite spectrum are norm dense in \(A_{\mathrm{sa}}\), and stable rank one means that invertible elements are norm dense in \(A\).

\textbf{Question.} Is it true that every unital separable complex \(C^*\)-algebra \(A\) satisfying at least one of
\[
\begin{aligned}
&\text{(i)}\quad A\cong A\otimes_{\min}\mathcal Z,\\
&\text{(ii)}\quad A\text{ is nuclear, has real rank zero, and has stable rank one}
\end{aligned}
\]
has \(\operatorname{gr}(A)\leq 1\)? Equivalently, under either set of hypotheses, is it true that for every \(a,b\in A_{\mathrm{sa}}\) and every \(\varepsilon>0\), there exist \(c,d\in A_{\mathrm{sa}}\) such that
\[
\|a-c\|<\varepsilon,\qquad \|b-d\|<\varepsilon,\qquad C^*(c,d)=A?
\]
No nuclearity assumption is added to (i), and no \(\mathcal Z\)-stability assumption is added to (ii).

The supplied record labels the problem partially solved. Its dated assessment reports dense single generators for important \(\mathcal Z\)-stable subclasses, including classifiable algebras with real rank zero and \(\mathcal Z\)-stable approximately subhomogeneous algebras, while reporting that the assertion in the stated generality remains open. This is supplied curation, not an independently verified mathematical or literature assessment; the precise hypotheses of the reported subclass results are not supplied. The report citation gives the year 2013, whereas the proposed year and DOI identifier indicate 2012; this bibliographic discrepancy is unresolved.
\subsection{Short English statement}
Must every unital separable \(C^*\)-algebra that either absorbs the Jiang--Su algebra or is nuclear with real rank zero and stable rank one have generating pairs of self-adjoint elements dense among all self-adjoint pairs? The supplied partially solved status and subclass results require independent review.
\subsection{Sources}
\begin{itemize}
\item Local UnsolvedMath registry, \texttt{data/unsolvedmath/problems.json}, permanent ID \(30002220\), source problem code \texttt{catalogue entry}, in the Oberwolfach Reports Open Problems collection. The supplied record gives proposed year 2012 and does not name a proposer.
\item Supplied report citation: \emph{C*-Algebras, Dynamics, and Classification} (2013), Oberwolfach Reports, pp.\ 3129--3209, DOI \texttt{10.4171/owr/2012/52}. \url{https://doi.org/10.4171/owr/2012/52}. The background describes Oberwolfach Reports as documenting mathematical workshops and their open-problem discussions.
\item Supplied literature reference, labelled ``generators in Z-stable real-rank-zero algebras'': \url{https://arxiv.org/abs/2006.08404}.
\item Supplied literature reference, labelled ``generator rank of subhomogeneous algebras'': \url{https://arxiv.org/abs/2006.03624}.
\item Supplied literature reference, labelled ``original OWR problem'': \url{https://doi.org/10.4171/owr/2012/52}.
\item Literature assessment date: \texttt{2026-08-21}. The partially solved label, the reported subclass results, and the claim that the general question remains open are supplied dated curation derived from the OWR collection, not independently verified here. The supplied background explicitly requests independent verification of time-sensitive status and cited claims.
\item The record labels its statement \texttt{corrected\_verified} and reports a manual semantic, typographic, and LaTeX review of a meaning-preserving cleanup. These are supplied provenance claims, not an independent verification of this draft or its sources.
\end{itemize}
\end{document}
